\section{Method}
\label{sec:method}

Figure~\ref{fig:architecture} shows where the library acts and how it learns; Figure~\ref{fig:usage} shows how
it decides which fields an agent uses. The rest of this section describes each step.

\begin{figure}[t]
\centering
% Where tracetrim acts. Schematic, no data. Colours as in the data figures.
\definecolor{ttblue}{HTML}{2A78D6}\definecolor{ttorange}{HTML}{EB6834}\definecolor{ttmuted}{HTML}{52514E}
\definecolor{ttpale}{HTML}{EAF2FC}\definecolor{ttpaleo}{HTML}{FDEDE6}\definecolor{ttline}{HTML}{BDBBB5}
\begin{tikzpicture}[font=\footnotesize, >=Stealth,
  col/.style={draw=ttline, rounded corners=3pt, fill=white},
  inner/.style={draw=ttline, rounded corners=2pt, fill=white, align=center, font=\scriptsize, minimum height=0.62cm},
  lbl/.style={font=\scriptsize, text=ttmuted, align=center},
  call/.style={->, line width=0.8pt},
  ret/.style={->, line width=0.8pt, draw=ttblue},
  lrn/.style={->, line width=0.8pt, draw=ttorange, dashed}]

% ===== (a) MCP ================================================================
\node[anchor=west, font=\footnotesize\bfseries] at (0,4.35) {(a) In front of the agent's MCP servers: no change to the agent's code};

% columns
\draw[col] (0,-0.5) rectangle (2.4,3.85);
\node[font=\footnotesize\bfseries] at (1.2,3.55) {Agent};
\node[lbl] at (1.2,1.5) {model and loop;\\sees the same\\tool names, plus\\\texttt{expand} and\\\texttt{tracetrim\_done}};

\draw[col, draw=ttblue, fill=ttpale] (5.0,-0.5) rectangle (11.0,3.85);
\node[font=\footnotesize\bfseries] at (8.0,3.55) {tracetrim proxy};

\draw[col] (13.0,-0.5) rectangle (15.4,3.85);
\node[font=\footnotesize\bfseries] at (14.2,3.55) {MCP servers};
\node[lbl] at (14.2,1.4) {billing, code,\\search\ldots\\[3pt]stdio or HTTP;\\headers, env\\or OAuth};

% inner boxes
\node[inner, draw=ttblue, minimum width=2.0cm] (proj) at (9.6,1.9) {projection};
\node[inner, minimum width=2.0cm] (orig) at (6.6,0.95) {originals};
\node[inner, draw=ttorange, fill=ttpaleo, minimum width=2.0cm] (lrn) at (9.6,0.05) {learner};

% 1. the call goes through unchanged
\draw[call] (2.4,3.0) -- node[lbl, above] {tool call} (5.0,3.0);
\draw[call] (5.0,3.0) -- node[lbl, below] {forwarded unchanged} (11.0,3.0);
\draw[call] (11.0,3.0) -- node[lbl, above] {same call} (13.0,3.0);
% 2. the result comes back, trimmed
\draw[call] (13.0,1.9) -- node[lbl, above] {full output} (proj.east);
\draw[ret] (proj.west) -- (5.0,1.9) -- node[lbl, above] {trimmed output\\+ \texttt{[trimmed]} line} (2.4,1.9);
\draw[->, ttline, line width=0.6pt] (proj.south west) ++(0.3,0) -- node[lbl, below right=-2pt] {stores} (orig.north east);
% 3. expand: answered from the stored original, no new server call
\draw[call] (2.4,1.1) -- node[lbl, above] {\texttt{expand}} (orig.west |- 0,1.1);
\draw[ret] (orig.west |- 0,0.8) -- node[lbl, below] {hidden fields} (2.4,0.8);
% 4. learning
\draw[lrn] (2.4,0.05) -- node[lbl, below] {\texttt{tracetrim\_done}\\(final answer)} (lrn.west);
\draw[lrn] (lrn.north) -- node[lbl, right] {reload} (proj.south);

% ===== learner, shared =========================================================
\node[draw=ttorange, fill=ttpaleo, rounded corners=3pt, text width=15.0cm, align=left, anchor=north west, inner sep=5pt,
      font=\scriptsize] (pipe) at (0,-0.85) {%
\textbf{Learner} (automatic mode). Each run is kept as a trace (calls, then the final answer) $\rightarrow$ usage
detection (Figure~\ref{fig:usage}) $\rightarrow$ telemetry report with suggestions, from the first run $\rightarrow$
after \texttt{min\_runs} runs, the automatic actions only: prudent projection, identifiers kept, \texttt{expand} on.
A tool whose outputs are expanded too often is no longer trimmed. Everything else stays in the report as a manual action.};

% ===== (b) function calling ====================================================
\node[anchor=west, font=\footnotesize\bfseries] at (0,-2.85) {(b) Function calling: three calls in the agent's own loop};
\node[inner, minimum width=2.4cm, minimum height=0.8cm, font=\footnotesize] (model) at (1.2,-3.65) {\textbf{Model}};
\node[inner, minimum width=2.9cm, minimum height=0.8cm] (loop) at (4.85,-3.65) {your loop runs\\the tool};
\node[inner, draw=ttblue, fill=ttpale, minimum width=3.6cm, minimum height=0.8cm] (auto) at (9.0,-3.65)
  {\texttt{AutoTrim.output()}\\\texttt{.expand()}\quad\texttt{.done(answer)}};
\node[inner, draw=ttorange, fill=ttpaleo, minimum width=2.4cm, minimum height=0.8cm] (l2) at (14.2,-3.65) {same learner};
\draw[call] (model.east) -- node[lbl, above] {tool call} (loop.west);
\draw[call] (loop.east) -- node[lbl, above] {full output} (auto.west);
\draw[lrn] (auto.east) -- (l2.west);
\draw[ret] (auto.south) -- ++(0,-0.35) -| node[lbl, below, pos=0.25] {trimmed output (no proxy, no \texttt{tracetrim\_done}: the loop has the answer)} (model.south);
\end{tikzpicture}

\caption{Where tracetrim acts. (a) As an MCP proxy, it stands in for the agent's servers: calls are forwarded
unchanged, only results are rewritten on their way back, and a trimmed result ends with a line listing what was
hidden, which \texttt{expand} returns from the stored original without calling the server again. The agent passes
its final answer to \texttt{tracetrim\_done}, which closes a run for the learner. (b) With function calling the same
learner sits in the agent's own loop. The automatic mode (dashed) is part of version 0.3 and was not evaluated in
Studies 1 and 2, which used offline learning and a fixed projection; Study 3 measures it through (b).}
\label{fig:architecture}
\end{figure}

\begin{figure}[t]
\centering
% How usage is detected, on an illustrative billing call (simplified values).
\definecolor{ttblue}{HTML}{2A78D6}\definecolor{ttorange}{HTML}{EB6834}\definecolor{ttmuted}{HTML}{52514E}
\definecolor{ttpale}{HTML}{EAF2FC}\definecolor{ttpaleo}{HTML}{FDEDE6}\definecolor{ttline}{HTML}{BDBBB5}
\begin{tikzpicture}[font=\footnotesize, >=Stealth,
  row/.style={anchor=west, inner sep=1pt, font=\ttfamily\scriptsize},
  tag/.style={anchor=west, inner sep=1pt, font=\scriptsize, text=ttmuted},
  hit/.style={->, line width=0.7pt, draw=ttblue}]

\node[anchor=west, font=\footnotesize\bfseries] at (0,4.0) {Tool output: one charge (simplified)};
\node[anchor=west, font=\footnotesize\bfseries] at (10.0,4.0) {What the agent writes after the call};

\draw[draw=ttline, rounded corners=3pt] (0,-0.75) rectangle (6.3,3.7);
\foreach \y/\t in {3.35/{id: "ch\_3QxB7"}, 2.85/{amount: 2900}, 2.35/{currency: "eur"},
                   1.85/{created: 1761897600}, 1.35/{status: "succeeded"},
                   0.85/{outcome.network\_status: "approved"}, 0.35/{billing\_details.address.* (6)},
                   -0.15/{payment\_method\_details.card.* (11)}}
  \node[row] at (0.15,\y) {\t};
\node[row, text=ttmuted] at (0.15,-0.5) {\ldots};

\node[tag, text=ttblue]   (t1) at (6.5,3.35) {used, as written};
\node[tag, text=ttblue]   (t2) at (6.5,2.85) {used, as \emph{29,00\,\texteuro} (S2)};
\node[tag]                     at (6.5,2.35) {never written};
\node[tag, text=ttblue]   (t4) at (6.5,1.85) {used, as a date (S2)};
\node[tag, text=ttorange]      at (6.5,1.35) {translated: \emph{not seen}};
\node[tag]                     at (6.5,0.85) {never written};
\node[tag]                     at (6.5,0.35) {never written};
\node[tag]                     at (6.5,-0.15) {never written};

\node[draw=ttline, rounded corners=3pt, fill=white, text width=5.0cm, align=left, anchor=north west, inner sep=5pt,
      font=\scriptsize] (ans) at (10.0,3.7) {%
  \emph{Final answer.} Two \emph{successful} charges of \textbf{29,00\,\texteuro} on \textbf{31 October 2025}, 47
  seconds apart (\textbf{ch\_3QxB7}, ch\_3QxC2): a double charge. Refund one of them.\\[4pt]
  \emph{Later call.} \texttt{list\_refunds(charge=}\\\texttt{"\textbf{ch\_3QxB7}")}};
\draw[hit] (t1.east) -- ($(ans.north west)+(0,-0.75)$);
\draw[hit] (t2.east) -- ($(ans.north west)+(0,-0.45)$);
\draw[hit] (t4.east) -- ($(ans.north west)+(0,-0.6)$);

\node[draw=ttblue, fill=ttpale, rounded corners=3pt, text width=15.0cm, align=left, anchor=north west, inner sep=5pt,
      font=\scriptsize] at (0,-1.05) {%
  \textbf{Matching.} A field is \emph{used} when one of its values appears, as a word $n$-gram ($n\le4$), in what the
  agent writes \emph{after} the call: later call arguments and the final answer. S2 also matches Unix times written as
  dates and minor units written as amounts.\\[2pt]
  \textbf{Projection} (prudent). Keep the fields used in any training run, the fields never judgeable, the identifying
  fields (S1: \texttt{id}, \texttt{*\_id}, \texttt{url}, \texttt{name}\ldots) and fields never seen; drop the rest. On a
  new run the dropped fields are listed in the \texttt{[trimmed]} line, and \texttt{expand} brings them back. A value
  the agent translates (\emph{succeeded} $\rightarrow$ \emph{successful}) is not matched: one of the failure
  mechanisms discussed in Section~\ref{sec:discussion}.};
\end{tikzpicture}

\caption{How usage is detected, on an illustrative billing call (values simplified). Blue: values matched in what
the agent wrote after the call, directly or through the safeguard S2; orange: a value the agent used after
translating it, which matching cannot see.}
\label{fig:usage}
\end{figure}

\paragraph{Traces.} A trace is the ordered record of what the agent was given and what it wrote: prompts, tool
calls (the arguments the agent wrote and the raw output it received) and messages, the final answer included. The
library reads a neutral JSONL format and the Python exports of the OpenAI Playground.

\paragraph{Usage.} Each structured output is flattened into leaves identified by a path in which list indices are
collapsed (\texttt{items[].keyword\_info.search\_volume}). A leaf is \emph{used} when one of its values appears in
something the agent wrote \emph{after} receiving it and not before. Values are matched as word $n$-grams
($n \le 4$): short strings as a whole, prose by any 4-gram, numbers in their plain form and with dot thousands separators, as Italian writes them (the library
matches comma separators too since version 0.3.4; the studies did not). Values too
generic to attribute (booleans, integers below 100, one-word strings shorter than four characters) are not
\emph{judgeable}. A value that already appears in the call's own arguments is an \emph{echo}: not a use, but the
record's key (the keyword next to its volume), and it is always kept.

\paragraph{Projection.} For each tool, the \emph{aggressive} projection keeps the used fields and the echoes; the
\emph{prudent} one also keeps every field that was never judgeable. Fields never seen while learning pass through,
so that a field added to an API is not hidden silently. Projected outputs are re-serialised without ASCII escapes.

\paragraph{Leave-one-out evaluation.} For each trace, the projection is learned from the other traces of the same
task and applied to it. Two quantities are reported: the tool-output tokens left (counted with the
\texttt{o200k\_base} tokenizer, each output once), and the \emph{misses}: calls in which the projection removes a
field that this trace used later. Usage is \emph{detected textual reuse}, neither a lower nor an upper bound of what
the agent uses: a field read only while reasoning leaves no trace in the text (reasoning is encrypted in the APIs
used here), and a value that coincides with another field's value, or a field kept for one element of a list, is
counted as used where it was not. ``After'' means a later model turn: calls emitted together in one turn are
generated before any of their outputs exist (deviation 14).

\paragraph{Learning source.} Usage can be learned from everything the agent wrote (\emph{full}) or only from the
arguments of later tool calls (\emph{calls}), which is what a proxy between the agent and its tools can observe.

\paragraph{Structural compression.} After projection, outputs can be cleaned (null and empty values removed) and
written in a TOON-style encoding \citep{schopplich2026toon} implemented in the library (not the reference
implementation): indented \texttt{key: value} lines, with arrays of uniform flat objects as tables whose field names
appear once.

\paragraph{\texttt{expand}.} At run time, a projected output is followed by a line listing the hidden fields,
grouped two levels deep and largest groups first, and the agent can retrieve any of them with an \texttt{expand}
tool. This turns a loss into a round trip, provided the agent knows what is missing.

\paragraph{Safeguards (version 0.2, Study 2).} S1 keeps identifying fields (\texttt{id}, \texttt{*\_id},
\texttt{doi}, \texttt{url}, \texttt{number}, \texttt{title}, \texttt{name} and similar) whatever the usage says.
S2 recognises transformed values: Unix timestamps and ISO dates also match as written dates, in Italian and English;
integers in money-named fields also match as major units with two decimals; URLs also match by their last segment.

\subsection{Detected reuse, required information and misses}
\label{sec:definitions}

Three terms are used throughout with a fixed meaning. \emph{Detected reuse}: a leaf path of a tool output is
detected as reused in a run when one of its judgeable values matches, by the rules above (and, with S2, the
transformations it recognises), something the agent wrote in a later model turn and not before. It is a property of
the trace and of the matcher. \emph{Required information}: for a task and a run, the leaf paths whose values the
agent needs to produce a correct answer, whether it writes them as they are, writes them transformed, or only reads
them to decide. It is defined by the task and its correct answers, not by the trace, and in general it cannot be
observed. \emph{Projection miss}: a held-out call in which the projection removes a path that the held-out run
reused, as detected. Misses are counted against detected reuse, not against required information: a projection
without misses has removed nothing the detector saw used, which is not the same as removing nothing required.

\paragraph{Why absence of detected reuse certifies nothing by itself.} A tool returns two issues, each with an
identifier and a title: issue A, titled ``staging deploy fails'', and issue B, titled ``flaky unit test''. In a first
task the agent must report the open issue with the earliest identifier; in a second, the issue that reports a failed
staging deployment; in both it answers with an identifier only. Issue A is both the earliest and the one about
staging, so the two tasks, though their instructions differ, can produce the same sequence of calls, outputs and answer,
and therefore the same detected-reuse signal, in every run. In neither is the title
ever detected as reused, since its value never appears in what the agent writes. A learner that relies on that signal cannot
tell the two tasks apart: if it removes the title in the first, where it is not needed, it removes it in the
second, where the agent needs it to choose. Study~1 met this case in CI triage (Section~\ref{sec:study1}). However
many traces are collected, the absence of detected reuse supports removing a field only under an added assumption,
for instance that every required value is written verbatim, which fields used only to reason violate. This is why a
learned projection is treated here as a proposal, with identifiers protected, the hidden fields listed and
\texttt{expand} as a way back.
